\providecommand{\BM}{\textbf{[B]}}
\providecommand{\KM}{\textbf{[K]}}
\providecommand{\SM}{\textbf{[S]}}
\providecommand{\XM}{\textbf{[X]}}
\providecommand{\QM}{\textbf{[Q]}}

\chapter*{Common Anchor: Bridge Equations to Observer Patch Holography and the Electroweak Scale as Meeting Point}
\addcontentsline{toc}{chapter}{Common Anchor: Bridge Equations to OPH and the Electroweak Scale as Meeting Point}

{\small\noindent\textbf{Abstract.}\quad\ignorespaces
Observer Patch Holography (OPH) lists as an open item that length and tick do
not follow from its own update rules. This document writes out the bridge to
FFGFT as equations. The OPH cell energy is exactly the FFGFT bit energy
$E_{\rm bit}=\hbar c/L$ at the cell edge $L_{\rm cell}=\sqrt P\,\ell_P$ \BM; the
collapse threshold of Doc.~329 gives $n_{\rm thr}=\sqrt{P/2}$ there, with the
boundary at $P=2$ \BM; $\tilde T\cdot m=1$ supplies a native tick
$\tilde T_{\rm cell}=\sqrt P\,t_P$. Below $v$, FFGFT delivers
all masses as ratios to $v$. The electroweak scale
$v$ is therefore the meeting point: the conditional OPH hierarchy $v/E_\star$
\QM{} and the FFGFT ladder $m_e/v$ \KM{} together give $m_e/E_P$ to
$0.93\,\%$. A single SI anchor then fixes both frameworks; through FFGFT's own chain the
cell edge and the cell tick become absolute. The coefficient of that reading
has meanwhile been derived on the OPH side (Doc.~376); \#740 itself remains
open there.
}\medskip

\section*{List of symbols}
\addcontentsline{toc}{section}{List of symbols}

\medskip
\begin{tabular}{@{}ll@{}}
\textbf{[B]} & algebraically proven \\
\textbf{[K]} & numerically confirmed \\
\textbf{[S]} & structurally plausible, derivation open \\
\textbf{[X]} & checked and negative \\
\textbf{[Q]} & taken from an external source, cited \\
\end{tabular}

\medskip
\begin{tabular}{@{}lp{10cm}@{}}
\toprule
\textbf{Symbol} & \textbf{Meaning} \\
\midrule
$\xi$ & FFGFT scaling parameter ($4/30000$) \\
$\ell_P$, $t_P$, $E_P$ & Planck length, time, energy (not reduced) \\
$L_0=\xi\ell_P$ & FFGFT resolution floor (Doc.~180) \\
$E_{\rm bit}(L)=\hbar c/L$ & system-dependent bit energy (Doc.~257) \\
$n_{\rm thr}(L)$ & collapse threshold, number of bit quanta up to Compton $=$ Schwarzschild (Doc.~329) \\
$M_{\rm coll}=m_P/\sqrt2$ & collapse mass (Doc.~329) \\
$\tilde T$ & intrinsic time scale, $\tilde T\cdot m=1$ \\
$v$ & electroweak scale ($246.22$\,GeV, from $G_F$) \\
$r_i$, $p_i$ & prefactor and exponent of the ladder $m_i=r_i\,\xi^{p_i}v$ \\
\midrule
$P$ & OPH pixel ratio $a_{\rm cell}/\ell_P^2$ \\
$L_{\rm cell}$, $E_{\rm cell}$ & OPH cell edge $\sqrt P\,\ell_P$, cell energy $E_P/\sqrt P$ \\
$E_\star$ & OPH reference energy of the hierarchy $v/E_\star$ \\
$\alpha_U$ & OPH unified coupling \\
$N$ & OPH capacity of the whole network \\
\bottomrule
\end{tabular}

\section{Occasion and classification}

Tracking item \#740 of the OPH repository records that length and tick do not
follow from the OPH updates: length is tied to an outside value only through
the capacity $N$ and $\Lambda\ell_P^2=3\pi/N$, and the tick only through
ratios. Doc.~364 examined the icosahedral OPH carrier spectrally; this
document treats the question of scale.

The roles are not symmetric. FFGFT derives from $\xi$; measured values and
Standard Model quantities serve only for conversion and comparison. OPH
quantities enter here as cited inputs from an external framework \QM{} and
carry the status their own documentation assigns them. The cosmic sector
stays bracketed (P39); $N$ is not used.

\section{Inputs}

\subsection*{FFGFT}
\begin{itemize}\itemsep2pt
\item Bit energy $E_{\rm bit}(L)=\hbar c/L$, system-dependent (Doc.~257).
\item Collapse threshold $n_{\rm thr}(L)=L/(\sqrt2\,\ell_P)$ and
  $M_{\rm coll}=m_P/\sqrt2$ from Compton $=$ Schwarzschild (Doc.~329) \BM.
\item Resolution floor $L_0=\xi\,\ell_P$ as a ratio (Docs.~180, 329).
\item $\tilde T\cdot m=1$: every energy has a proper time $\hbar/E$.
\item Lepton ladder $m_e=\tfrac43\,\xi^{3/2}\,v$ (Docs.~006, 373) \KM.
\end{itemize}

\subsection*{OPH \QM}
\begin{itemize}\itemsep2pt
\item Pixel ratio $P=a_{\rm cell}/\ell_P^2$ as the fixed point of
  $P=\varphi+\sqrt\pi/A_T(P)$; public root $P=1.630968\ldots$, source root
  $1.630972\ldots$; cell energy $E_{\rm cell}=E_P/\sqrt P$
  (\texttt{code/P\_derivation/FULL\_DERIVATION.md}). The public root contains
  an empirical hadronic contribution; the source root does not.
\item Hierarchy $v/E_\star=P^{-1/2}\exp[-2\pi/(4\alpha_U(P))]
  =2.01998\times10^{-17}$ (Theorem~A,
  \texttt{code/particles/hierarchy/theorem\_package.md}); explicitly
  conditional, the physical identity of $E_\star$ is open in OPH.
\end{itemize}

\section{Bridge equations for length and tick}\label{sec:bruecke}

All statements in this section are dimensionless in Planck units; no SI anchor
is consumed.

\paragraph{B1 --- dictionary \BM.} With $L_{\rm cell}=\sqrt P\,\ell_P=1.2771\,\ell_P$,
\begin{equation}
  E_{\rm cell}^{\rm OPH}=\frac{E_P}{\sqrt P}=\frac{\hbar c}{L_{\rm cell}}=E_{\rm bit}(L_{\rm cell}).
\end{equation}
Both frameworks use the same energy--length relation. OPH fixes the length by
its fixed point; in FFGFT it stays system-dependent.

\paragraph{B2 --- collapse threshold \BM.} Inserted into Doc.~329:
\begin{equation}
  n_{\rm thr}(L_{\rm cell})=\sqrt{P/2}=0.9030,\qquad
  \frac{r_s(E_{\rm cell})}{L_{\rm cell}}=\frac2P=1.2263 .
\end{equation}
The boundary above which a single cell quantum stays outside its own
Schwarzschild radius is $P=2$. The OPH root lies below it: a cell quantum
localised on one cell sits at the collapse scale. For a holographic pixel this
is expected rather than objectionable, but it is a sharp condition that any
future change of $P$ must respect or explain. Source and public roots differ
here by $1.2\times10^{-6}$; the statement does not depend on the open hadronic
part.

\paragraph{B3 --- tick.} $\tilde T\cdot m=1$ assigns the cell a proper time:
\begin{equation}
  \tilde T_{\rm cell}=\frac{\hbar}{E_{\rm cell}}=\sqrt P\,t_P=\frac{L_{\rm cell}}{c}.
\end{equation}
As arithmetic this is \BM; physically it gives OPH a tick per cell in the same
unit as its length, without a master clock. Tying it to a laboratory clock
is possible but not yet carried out on the OPH side \SM. A laboratory clock,
such as the caesium transition, is itself a matter system with energy
$\Delta E_{\rm clk}$ and, by $\tilde T\cdot m=1$, its own tick $\hbar/\Delta E_{\rm clk}$.
The attachment requires the ratio $\Delta E_{\rm clk}/E_{\rm cell}$, computed within
OPH itself; OPH lists this step as open (clock frontier
$\varepsilon_{\rm clk}=\Delta E_{\rm clk}/E_\star$, caesium route as an incomplete
candidate). In FFGFT the question is not open: every mode has its tick
$\tilde T=1/m$, and tick ratios are mass ratios. FFGFT thus supplies the
translation principle, not the missing OPH value.

\paragraph{B4 --- floor \KM.} $L_{\rm cell}/L_0=\sqrt P/\xi=9578$. The OPH
cell lies about four orders of magnitude above the FFGFT resolution floor; no
conflict.

\section{The meeting point $v$}\label{sec:treffpunkt}

The two frameworks cover different sections of the scale ladder. OPH computes,
conditionally, from the top down to $v$; FFGFT from $v$ downwards: all
charged lepton masses via the ladder $m_i=r_i\,\xi^{p_i}v$ (Docs.~006, 373),
the boson masses via $\sin^2\theta_W=2/9$ and the trace rule (Docs.~372, 373).

\begin{center}
\begin{tabular}{@{}llr@{}}
\toprule
\textbf{Quantity} & \textbf{Origin} & \textbf{vs.\ measured} \\
\midrule
$v/E_\star=2.01998\times10^{-17}$ & OPH, Theorem A \QM & $+0.16\,\%$ (with $E_\star=E_P$) \\
$m_e/v=\tfrac43\,\xi^{3/2}=2.0528\times10^{-6}$ & FFGFT, Docs.~006/373 \KM & $-1.09\,\%$ (bare) \\
$m_e/E_P=4.1466\times10^{-23}$ & product & $-0.93\,\%$ \\
\bottomrule
\end{tabular}
\end{center}

For the reference energy $E_\star$ see the following subsection. The $-1.09\,\%$ on the FFGFT side is the known gap between
bare and measured lepton masses (Doc.~352), not a new finding. The product
depends on the OPH branch only at the $10^{-4}$ level.

Consequence: a single SI anchor fixes both frameworks. With $G$ as anchor, the
OPH hierarchy leads to $v$, and the FFGFT ladder leads from there to all
particle masses; with $m_e$ as anchor, the chain runs the other way. The
residual of about $1\,\%$ is the common consistency test. FFGFT also reaches
$m_e/E_P$ through its own chain; the meeting point $v$ is thus a comparison
of two independent routes.

\subsection*{The reference energy $E_\star$}\label{sec:estar}

\paragraph{What $E_\star$ is.} In OPH, $E_\star$ is the energy unit of the
calculation. The code sets the Planck energy equal to one; the cell energy is
$E_P/\sqrt P$, the unification scale $M_U=E_P\,e^{-2\pi}P^{1/6}$, the hierarchy
$v/E_\star=P^{-1/2}\exp[-2\pi/(4\alpha_U)]$ (reproduced, check script D1). Which
physical energy this unit is, OPH explicitly leaves open; it is to be fixed
through a clock, $E_\star=h\nu_{\rm clk}/\varepsilon_{\rm clk}$, that is, through the
same open step as in B3.

\paragraph{Two candidates.} The Planck energy exists in two conventions:
non-reduced $E_P=\sqrt{\hbar c^5/G}=1.2209\times10^{19}$\,GeV and reduced
$\bar E_P=E_P/\sqrt{8\pi}=2.435\times10^{18}$\,GeV. They differ by
$\sqrt{8\pi}\approx5.01$, depending on whether the gravitational coupling is
written as $G$ or as $8\pi G$. With measured $v$, OPH gives
$E_\star=1.2189\times10^{19}$\,GeV: $0.16\,\%$ below $E_P$, but a factor of
$5.0$ above $\bar E_P$ \KM. The reduced reading is thereby excluded; the
non-reduced one fits.

\paragraph{How sharp $0.16\,\%$ is.} $\alpha_U$ sits in the exponent. The
sensitivity is
\begin{equation}
  \frac{\partial\ln(v/E_\star)}{\partial\alpha_U}=\frac{\pi}{2\alpha_U^2}\approx929 .
\end{equation}
The deviation of $0.16\,\%$ corresponds to a shift of $\alpha_U$ by
$1.7\times10^{-6}$, relatively $4\times10^{-5}$ \BM. The same sensitivity
accounts exactly for the difference between the public and the source branch
($-8.4\times10^{-5}$, check script D5). The agreement with $E_P$ is therefore
not a rough match of orders of magnitude but a tight condition on $\alpha_U$.

\paragraph{Why a reading \SM.} OPH itself does not claim $E_\star=E_P$; the
assignment here follows from the numerical comparison. It would be confirmed
if OPH closes the clock attachment and obtains $E_\star=E_P$, or if OPH
singles out the coupling without $8\pi$ structurally. For the comparison in
§\ref{sec:wege} the reading is not critical: it only fixes against which of
the two Planck energies the $0.16\,\%$ is measured.

\section{Two routes to $m_e/E_P$ and their accuracy}\label{sec:wege}

FFGFT reaches the Planck scale through its own chain (Docs.~012, 013, 180):

\begin{enumerate}\itemsep2pt
\item $G=\dfrac{\xi^2}{4m_e}\,C_{\rm conv}\,K_{\rm frak}=6.6745\times10^{-11}\,
  \mathrm{m^3kg^{-1}s^{-2}}$, $+0.003\,\%$ against CODATA. The notation in
  Doc.~180 splits the same value into $C_{\rm dim}\cdot C_{\rm conv}$; both
  forms are equivalent.
\item $\ell_P=\sqrt{\hbar G/c^3}=1.6163\times10^{-35}$\,m and hence
  $L_0=\xi\,\ell_P=2.155\times10^{-39}$\,m.
\item $E_P=1.22087\times10^{22}$\,MeV, $m_e/E_P=4.1855\times10^{-23}$
  ($+0.002\,\%$) \KM.
\item $v/E_P=2.01676\times10^{-17}$.
\end{enumerate}

Not every route of calculation yields the same accuracy:

\begin{center}
\small
\begin{tabular}{@{}llrl@{}}
\toprule
\textbf{Route} & \textbf{Quantity} & \textbf{Dev.} & \textbf{Origin of the deviation} \\
\midrule
FFGFT chain $\xi\to G\to E_P$ & $m_e/E_P$ & $+0.002\,\%$ & accuracy of $G$ \\
FFGFT ladder & $m_e/v$ & $-1.09\,\%$ & bare versus measured (Doc.~352) \\
OPH, Theorem A \QM & $v/E_\star$ vs.\ FFGFT chain & $+0.16\,\%$ & reading $E_\star=E_P$, conditional fixed point \\
Combination via $v$ & $m_e/E_P$ & $-0.93\,\%$ & sum of the two preceding \\
\bottomrule
\end{tabular}
\end{center}

The routes do not contradict each other. Each deviation is assigned to a
known place: the bare residual of the lepton ladder on one side, the
identification of $E_\star$ on the other. For the comparison table, the route
by which a value was obtained must therefore always be stated alongside the
value.

\paragraph{Consequence for \#740.} With FFGFT's own chain the bridge
equations become absolute. The OPH cell edge is
$L_{\rm cell}=\sqrt P\,\ell_P=2.064\times10^{-35}$\,m, the cell tick
$\tilde T_{\rm cell}=\sqrt P\,t_P=6.885\times10^{-44}$\,s \KM; every further
tick ratio is, by $\tilde T\cdot m=1$, a mass ratio. Length and tick thus
follow without the detour via $N$ and $\Lambda$, under the reading
$E_\star=E_P$. Its coefficient has meanwhile been derived on the OPH side
(Doc.~376): from two observers seeing different stress it follows that
$\kappa=b\,L^2/q$ and hence $G=L^2$ \QM. \#740 itself remains open there;
the length identity, the clock and a matching factor are still carried as
open. The question of scale within \#740 is thus answered on the FFGFT side,
the item as a whole is not.

\section{What the bridge does not achieve}

\begin{itemize}\itemsep2pt
\item The upper part of the chain comes from an external framework and is
  conditional there \QM; the combination is a consistency test between the
  frameworks.
\item Unit of time: in FFGFT every time scale follows from $\tilde T=1/m$; the
  SI second serves only for conversion. In OPH the unit of time is tied to
  $E_\star$ and is to be fixed through the clock attachment; this step is open
  on the OPH side \SM. Through the FFGFT chain the cell tick is already
  absolute under the reading $E_\star=E_P$; the laboratory clock is then a check.
\item $N$ and the cosmic sector are not involved (P39).
\item The algebraic carriers still differ ($A_5$ versus $T^4/\mathbb Z_3$,
  Doc.~364); the bridge concerns scales, not symmetries.
\end{itemize}

\section*{Tool and method}
\addcontentsline{toc}{section}{Tool and method}
Claude (Anthropic) was used as a tool for: reading the OPH source files,
writing and running the Python check script, drafting the \LaTeX{} sources,
and maintaining register and changelog. The tool computes and writes; the
author decides what is computed and written, judges every result and
determines what it means physically. Every statement is secured by assertions
in the check script. OPH values are taken verbatim from the public
repository; measured values (CODATA 2022, PDG) serve only as comparators.

\section{Result}

\begin{enumerate}
\item B1: $E_{\rm cell}^{\rm OPH}=E_{\rm bit}(\sqrt P\,\ell_P)$ \BM.
\item B2: $n_{\rm thr}(L_{\rm cell})=\sqrt{P/2}$, boundary $P=2$, OPH below it \BM.
\item B3: cell tick $\tilde T_{\rm cell}=\sqrt P\,t_P$ from $\tilde T\cdot m=1$;
  laboratory attachment open on the OPH side \SM.
\item B4: $L_{\rm cell}/L_0=\sqrt P/\xi\approx 9578$ \KM.
\item Meeting point $v$: OPH $v/E_\star$ \QM{} $\times$ FFGFT $m_e/v$ \KM{}
  gives $m_e/E_P$ to $-0.93\,\%$; a single SI anchor suffices for both
  frameworks.
\item Length and tick absolute through the FFGFT chain: $L_{\rm cell}=2.064\times10^{-35}$\,m,
  $\tilde T_{\rm cell}=6.885\times10^{-44}$\,s \KM; the coefficient of the reading
  derived on the OPH side, \#740 as a whole open there (Doc.~376).
\item FFGFT's own chain $\xi\to G\to E_P$: $m_e/E_P$ to $+0.002\,\%$ \KM;
  the accuracy depends on the route of calculation (§\ref{sec:wege}).
\end{enumerate}

\medskip
\noindent\textbf{Check script:} \texttt{2/python/Dok374\_Skripte/pruef\_374\_anker.py} (30/30)

\medskip
\noindent\textbf{Register (Doc.~190):} no entry.

\begin{thebibliography}{9}
\bibitem{dok006} J.~Pascher, \emph{Doc.~006 --- T0 particle masses} (2025).
\bibitem{dok012} J.~Pascher, \emph{Doc.~012 --- T0 gravitational constant} (2025).
\bibitem{dok013} J.~Pascher, \emph{Doc.~013 --- T0 and SI units} (2025).
\bibitem{dok180} J.~Pascher, \emph{Doc.~180 --- Derivation of $L_0$} (2026).
\bibitem{dok257} J.~Pascher, \emph{Doc.~257 --- Bit energy} (2026).
\bibitem{dok329} J.~Pascher, \emph{Doc.~329 --- $n_{\rm thr}$ scale analysis} (2026).
\bibitem{dok352} J.~Pascher, \emph{Doc.~352 --- Lepton residuals} (2026).
\bibitem{dok364} J.~Pascher, \emph{Doc.~364 --- The icosahedral carrier and the $D_4$ lattice} (2026).
\bibitem{dok373} J.~Pascher, \emph{Doc.~373 --- The Yukawa mechanism} (2026).
\bibitem{oph} B.~Mueller et al., \emph{Observer Patch Holography},
  \texttt{github.com/FloatingPragma/observer-patch-holography} (as of 22~September~2026).
\end{thebibliography}
